\vspace{0.5em}\noindent\textbf{RECOMMENDATION SYSTEMS FAIL: DATA, NOT
ALGORITHM}\par

When building a product recommendation system with Amazon Personalize
during an internship at FCAJ, we discovered that poor performance
usually comes from data quality, not algorithm choice.

\vspace{0.5em}\noindent\textbf{3 Key Lessons}\par

\vspace{0.5em}\noindent\textbf{1. Understand Your Algorithm's Core
Mechanism}\par

\begin{itemize}
\tightlist
\item
  Amazon Personalize uses \textbf{collaborative filtering} --- it finds
  users with similar behavior and recommends based on those patterns.
\item
  \textbf{Key insight:} If user interactions are random, no patterns
  exist for the algorithm to learn.
\end{itemize}

\subsubsection{2. Data Quality \textgreater{} Algorithm
Selection}

\begin{itemize}
\tightlist
\item
  Changing recipes without fixing underlying data issues wastes time and
  money.
\item
  \textbf{Focus on data structure:}

  \begin{itemize}
  \tightlist
  \item
    Create realistic user segments (e.g., electronics lovers, fashion
    fans, bookworms)
  \item
    Simulate real behavior patterns (sessions, conversion funnels,
    power-law distribution)
  \item
    The algorithm can only learn what your data reveals
  \end{itemize}
\end{itemize}

\vspace{0.5em}\noindent\textbf{3. Synthetic Data Must Reflect
Reality}\par

\begin{itemize}
\tightlist
\item
  \texttt{random.choice()} is fast but useless for meaningful testing.
\item
  \textbf{Design synthetic data carefully:}

  \begin{itemize}
  \tightlist
  \item
    Model user sessions with related product views
  \item
    Apply conversion rates (view\(\rightarrow\)cart,
    cart\(\rightarrow\)purchase)
  \item
    Build distinct user segments with clear preferences
  \end{itemize}
\end{itemize}

\vspace{0.5em}\noindent\textbf{Results Comparison: Same Algorithm, Better
Data}\par

{\def\LTcaptype{none} % do not increment counter
\begingroup
\small
\setlength{\tabcolsep}{3pt}
\renewcommand{\arraystretch}{1.15}
\begin{longtable}[]{@{}llll@{}}
\toprule\noalign{}
Metric & Random Data & Structured Data & Improvement \\
\midrule\noalign{}
\endhead
\bottomrule\noalign{}
\endlastfoot
\textbf{Precision@5} & 0.0889 & 0.4348 & \textbf{4.9x} \\
\textbf{NDCG@10} & 0.1799 & 0.6512 & \textbf{3.6x} \\
\textbf{MRR@25} & 0.1216 & 0.7130 & \textbf{5.9x} \\
\textbf{Coverage} & 0.8218 & 0.9505 & \textbf{+15.7\%} \\
\end{longtable}
\endgroup
}

\vspace{0.5em}\noindent\textbf{Key Takeaways:}\par

\begin{itemize}
\tightlist
\item
  MRR increased nearly 6x --- recommendations now appear at the top of
  the list where users actually see them
\item
  Coverage improved by 15.7\% --- the system recommends more products,
  not just best-sellers
\item
  A well-structured dataset can multiply performance 5-6 times without
  changing any algorithm
\end{itemize}

\vspace{0.5em}\noindent\textbf{Component Breakdown}\par

\textbf{Data Preparation:} The critical foundation --- includes user
segmentation, behavior modeling, and realistic interaction patterns.

\textbf{Amazon Personalize:} Managed service with pre-built recipes (AWS
provides the algorithms, you provide the data).

\textbf{Evaluation Metrics:} Precision@5, NDCG@10, MRR@25, Coverage ---
measure recommendation quality and diversity.

\textbf{Structured Data:} Requires user segments (electronics, fashion,
books), session-based interactions, conversion funnels, and power-law
distributions.

\vspace{0.5em}\noindent\textbf{Key Takeaway}\par

Tools and algorithms are only 20\%. The remaining 80\% is \textbf{Data
Quality} and \textbf{Understanding Your Problem Domain}.

When metrics are abnormally low, \textbf{examine your data before
changing the algorithm}.

\begin{center}\rule{0.5\linewidth}{0.5pt}\end{center}

\textbf{Original post:}
\href{https://www.facebook.com/groups/awsstudygroupfcj/permalink/2229279517837008/?rdid=7zgI0zX0X33fQod5\#}{Facebook
AWS Study Group FCJ}
